\documentclass[11pt,a4paper,onecolumn,final]{article}

\usepackage[T1]{fontenc}
\usepackage{graphicx}
\usepackage{float}
\usepackage{chngcntr}
\counterwithin{figure}{section}

\usepackage{tikz}
\usetikzlibrary{angles, quotes}

\usepackage{caption}
%Input encoding for typing directly Greek characters
%Other values for the option: latin1, ascii, ansinew
%It is generally advised to put it before babel
\usepackage[utf8]{inputenc}

%MAIN Language the LAST one
\usepackage[greek,english]{babel}


\usepackage[
backend=biber,
style=alphabetic,
sorting=ynt
]{biblatex}
\addbibresource{sample.bib}

%\usepackage{kerkis}

%The teubner package extends the facilities offered by the CB Greek fonts
%NOTE: This includes also the graphicx package!
%NOTE: When the teubner package is loaded, command \cap (denoting set intersection) is overwritten
%NOTE: TEUBNER.STY FILE HAS BEEN CHANGED A BIT IN LINES 255-257 WITH 3 REPLACEMENTS \cap -> \capt
\usepackage{teubner}

%Package for displaying math
\usepackage{amsmath}

%Loading a superset of amsfonts
\usepackage{amssymb}

%Provides an enhanced version of \newtheorem
%1.Defines * for unnumbered environments
%2.Defines a proof environment
%3.Defines plain, definition, and remark theorem styles
\usepackage{amsthm}

%Allows text in eps files to be replaced with LATEX symbols
\usepackage{psfrag}

\makeatletter\let\c@lofdepth\relax\let\c@lotdepth\relax\makeatother
\usepackage[small,bf,tight]{subfigure}

%Math symbols of the dsfont style. Enables the use of \mathds{}
\usepackage{dsfont}

%Defines several symbols in text mode (e.g. the \textreferencemark for use in the itemize environment)
\usepackage{textcomp}

%Enables the Ralph Smith's Formal Script font in math mode via the \mathscr{} command
\usepackage{mathrsfs}

\usepackage{geometry}
%PDF VIEW
\geometry{total={210mm,297mm},left=25mm,right=25mm,bindingoffset=0mm, top=25mm,bottom=25mm}
%PRINT
%\geometry{total={210mm,297mm},left=20mm,right=20mm,bindingoffset=10mm, top=25mm,bottom=25mm}

\usepackage[breaklinks=true,colorlinks=true,
%linkcolor=blue,urlcolor=blue,citecolor=blue,% PDF VIEW
linkcolor=black,urlcolor=black,citecolor=black,% PRINT
bookmarks=false,bookmarksopenlevel=2]{hyperref}
\usepackage{hyperref}

\usepackage{cleveref}

%Package to use align justify
\usepackage{ragged2e}

%Convenient language switch
\newcommand{\gr}[1]{\foreignlanguage{greek}{#1}}
\newcommand{\en}[1]{\foreignlanguage{english}{#1}}

\begin{document}

%COVER PAGE%%%%%%%%%%%%%%%%%%%%%%%%%%%%%%%%%%%%%%%%%%%%%%%%%%%%%%%%%%%%%%%%%%%%%
\thispagestyle{empty}
{\sffamily\centering\Large

{\Large \textbf{Τεχνικές βελτιστοποίησης} \\[0.5cm]

\vspace{1.0cm}

{\Large \textbf{Project}\\
\textbf{Ελαχιστοποίηση συνάρτησης πολλών μεταβλητών} \\
\textbf{Γενετικοί αλγόριθμοι} \\[0.5cm]

\vspace{1.0cm}

{\LARGE Λαβδάκης Γεώργιος \\
ΑΕΜ:10672}

\vspace{1.0cm}

\newpage

\tableofcontents

\newpage

\begin{flushleft}

\section{Εισαγωγή}

Το project αφορά την μελέτη γενετικών αλγορίθμων. Οι γενετικοί αλγοριθμοι στοχεύουν στην εύρεση βέλτιστων λύσεων σε πολύπλοκα προβλήματα τα οποία έχουν πολλές παραμέτρους και άρα δεν μπορούν να εφαρμοστούν εύκολα οι κλασσικές μέθοδοι βελτιστοποίησης. Αποτελούν τεχνικές ολικής βελτιστοποίησης και είναι εμπνευσμένοι από τη βιολογία. Συγκεκριμένα, χρησιμοποιούν την ιδέα της εξέλιξης μέσω γενετικής μετάλλαξης, φυσικής επιλογής και διασταύρωσης, επιλέγοντας τα άτομα με την καλύτερη απόδοση ώστε να μεταβιβάσουν τα χαρακτηριστικά τους στους απογόνους.  Έτσι, με το πέρασμα των γενεών ο πληθυσμός των ατόμων εξελίσσεται προς τα άτομα με την καλύτερη απόδοση και η διαδικασία αυτή επαναλαμβάνεται μέχρι να εντοπιστεί η βέλτιστη λύση.

\vspace{0.4cm}

Οι γενετικοί αλγόριθμοι ξεκινάνε με ένα αρχικό πλήθος από τυχαίες λύσεις του προβλήματος, που ικανοποιούν τις προδιαγραφές του και ονομάζεται πληθυσμός (population). Κάθε άτομο του πληθυσμού αντιπροσωπεύει ένα διάνυσμα λύσης του προβλήματος το οποίο ονομάζεται χρωμόσωμα (chromosome). Τα χρωμοσώματα αυτά, όπως και στη φύση, εξελίσσονται μέσω επαναλήψεων που ονομάζονται γενεές (generations) και κάθε γενιά αξιολογείται με βάση μια συνάρτηση ικανότητας (fitness function). Για τη δημιουργία κάθε επόμενης γενιάς, των απογόνων (offspring), χρησιμοποιούνται οι πράξεις της  διασταύρωσης (crossover), μετάλλαξης (mutation) και επιλογής (selection) με βάση τη fitness function. Μετά από αρκετές γενεές, ο αλγόριθμος εντοπίζει το καλύτερο χρωμόσωμα, το οποίο αποτελεί την βέλτιστη ή υποβέλτιστη λύση στο πρόβλημα, σύμφωνα με τη συνθήκη τερματισμού \(|\frac{f(x_{k-1})-f(x_{k})}{f(x_{k-1})}|<\epsilon\) με \(\epsilon\) την επιθυμητή ακρίβεια της λύσης. Παρακάτω φαίνεται η δομή ενός γενετικού αλγορίθμου.

\begin{figure}[H]
    \centering
    \includegraphics[width=0.7\linewidth]{mceclip0_1609571964.png}
    \caption{Δομή ενός γενετικού αλγορίθμου}
    \label{fig:enter-label}
\end{figure}

\newpage

\section{Θέμα 1: Μαθηματική διατύπωση προβλήματος}

Στόχος είναι η ελαχιστοποίηση του συνολικού χρόνου διάσχισης του δικτύου ανά όχημα ως προς τον ρυθμό διέλευσης οχημάτων στον δρόμο \(i \in [1, 17]\). Ο όρος \(x_{i} \cdot T_{i}\) εκφράζει τον συνολικό χρόνο που ξοδεύουν ανά λεπτό τα οχήματα του δρόμου i, άρα διαιρώντας το άθροισμά τους με τον ρυθμό εισερχόμενων οχημάτων V προκύπτει ο χρόνος ανά όχημα. Επομένως, επιλέγεται συνάρτηση προς ελαχιστοποίηση η

\[f(x) = \frac{1}{V}\sum_{i=1}^{17} x_{i} \cdot T_{i}\]

όπου \(x_{i}\) [οχ./min] ο ρυθμός διέλευσης οχημάτων στο δρόμο i και \(T_{i}\) ο χρόνος κίνησης στον ίδιο δρόμο. Ο \(T_{i}\) συναρτήσει του \(x_{i}\) είναι

\[T_{i}(x_{i})=t_{i}+a_{i}\frac{x_{i}}{1-\frac{x_{i}}{c_{i}}} [min]\] 

όπου \(a_{i}\) σταθερές, \(c_{i}\) [οχ./min] ο μέγιστος δυνατός ρυθμός διέλευσης οχημάτων από τον δρόμο i, \(t_{i}\) [min] ο σταθερός χρόνος που απαιτείται για να κινηθούμε στον ίδιο δρόμο όταν η κίνηση είναι ασθενής και \(V = 100\) [οχ./min] ο ρυθμός εισερχόμενων οχημάτων στο δίκτυο. Οι τιμές των \(t_{i}\) δεν δίνονται στην εκφώνηση, γι' αυτό θεωρούμε \(t_{i}=0\). Σημειώνεται ότι ο όρος \(\sum_{i} x_{i} t_{i}\) γενικά εξαρτάται από τα \(x_{i}\), άρα αν οι τιμές των \(t_{i}\) ήταν γνωστές θα έπρεπε να συμπεριληφθεί στη συνάρτηση. Επίσης, επειδή οι γενετικοί αλγόριθμοι είναι αλγόριθμοι μεγιστοποίησης τελικά ως fitness function ορίζεται η \(\frac{1}{f(x)}\). 

\vspace{0.4cm}

Για να αποφευχθεί η συγκέντρωση οχημάτων στους κόμβους του δικτύου είναι επιθυμητό όσα οχήματα εισέρχονται σε κάθε κόμβο τόσα και να εξέρχονται. Επιπλέον, ο \(x_{i}\) δεν πρέπει να ξεπερνάει τον \(c_{i}\) γιατί τότε ο χρόνος διάσχισης της διαδρομής γίνεται άπειρος. Σύμφωνα με αυτές τις δύο συνθήκες και με βάση τη μορφή του δικτύου, προκύπτουν οι παρακάτω περιορισμοί  

\(\cdot x_{1}+x_{2}+x_{3}+x_{4}=V=100\)

\(\cdot x_{1}-x_{5}-x_{6}= 0\)

\(\cdot x_{2}-x_{7}-x_{8}=0\)

\(\cdot x_{3}+x_{8}+x_{9}-x_{11}-x_{12}-x_{13}=0\)

\(\cdot x_{4}-x_{9}-x_{10}=0\)

\(\cdot x_{6}+x_{7}+x_{13}-x_{14}-x_{15}=0\)

\(\cdot x_{10}+x_{11}-x_{17}=0\)

\(\cdot x_{5}+x_{14}-x_{16}=0\)

\(\cdot x_{12}+x_{15}+x_{16}+x_{17}=V=100\)

\(\cdot 0 \leq x_{i} < c_{i}, \; i \in [1, 17]\)

\newpage

\section{Θέμα 2: Υλοποίηση γενετικού αλγορίθμου}

Παρακάτω παρουσιάζεται ένας γενετικός αλγόριθμος στο Matlab που να λύνει το πρόβλημα. Επιλέγεται κωδικοποίηση με πραγματικές τιμές (real coded genetic algorithm) και κάθε χρωμόσωμα να είναι ένα διάνυσμα 1x17 που περιέχει όλα τα \(x_{i}, i \in [1, 17]\). Το αρχείο στο οποίο υλοποιείται ο αλγόριθμος είναι το \textbf{GeneticAlgorithm.m}.  

\vspace{0.4cm}

\(\cdot\) Αρχικά, δημιουργείται η πρώτη γενιά με τη συνάρτηση \textbf{firstGeneration.m}. Κάθε επόμενη γενιά δημιουργείται με τη συνάρτηση \textbf{nextGeneration.m}.

\vspace{0.4cm}

\(\cdot\) Ο έλεγχος των περιορισμών γίνεται με τη συνάρτηση \textbf{constraints.m}, η οποία επιστρέφει 0 αν κάποιος δεν ικανοποιούνται και 1 αν ικανοποιούνται όλοι. Για τον έλεγχο των ισοτικών περιορισμών π.χ. \(x_{1}+x_{2}+x_{3}+x_{4}=100\), ελέγχονται οι ανισότητες π.χ. \(|x_{1}+x_{2}+x_{3}+x_{4}-100| < \epsilon\) όπου \(\epsilon = 10^{-8}\) το μέγιστο ανεκτό σφάλμα λόγω της αριθμητικής κινητής υποδιαστολής. Ελέγχεται επίσης ότι για κάθε δρόμο ισχύει \(0 \leq x_{i} < c_{i}\).

\vspace{0.4cm}

\(\cdot\) Μετά, η fitness function, στο αρχείο \textbf{fitnessFunction.m}, αξιολογεί κάθε χρωμόσωμα του πληθυσμού επιστρέφοντας μία βαθμολογία για το καθένα (fitness value) η οποία αντιπροσωπεύει την αποτελεσματικότητα της λύσης. Όσο μεγαλύτερη η τιμή που επιστρέφει η fitness function τόσο καλύτερο είναι και το αποτέλεσμα, δηλαδή τόσο πιο κοντά είναι η λύση στην βέλτιστη. Ο χρόνος διάσχισης ανά όχημα \(f(x)\) υπολογίζεται στη συνάρτηση \textbf{totalTime.m}.

\vspace{0.4cm}

\(\cdot\) Καθώς η διαδικασία επιλογής στους γενετικούς αλγορίθμους γίνεται στοχαστικά, κανένα χρωμόσωμα στον πληθυσμό, όσο ικανό και να είναι, δεν περνάει στην επόμενη γενιά με βεβαιότητα. Γι΄ αυτό, η διαδικασία επιλογής γίνεται με τον τροχό της ρουλέτας. Η επιλογή με ρουλέτα, στο αρχείο \textbf{rouletteSelection.m}, αρχικά υπολογίζει το συνολικό fitness του πληθυσμού, βρίσκει τη πιθανότητα επιβίωσης κάθε χρωμοσώματος, η οποία είναι ανάλογη του fitness value του \(P_{i}=\frac{fitness \: value_{i}}{total \: fitness}\) και υπολογίζει την αθροιστική πιθανότητά του. Μετά, αφού ταξινομήσει τις πιθανότητες σε αύξουσα σειρά, παράγει τυχαίο αριθμό στο διάστημα \([0, 1]\) και επιλέγεται το πρώτο χρωμόσωμα με αθροιστική πιθανότητα μεγαλύτερη από τον τυχαίο αριθμό. Η διαδικασία επαναλαμβάνεται 10 φορές, μία για κάθε θέση του πληθυσμού, και έτσι επιλέγεται η επόμενη γενιά. Ουσιαστικά μοιράζεται ο τροχός της ρουλέτας σε τμήματα και αν ένα χρωμόσωμα έχει μεγάλο fitness value, καταλαμβάνει μεγαλύτερο μέρος του κύκλου και άρα έχει περισσότερες πιθανότητες να επιλεχθεί ενώ αν ένα χρωμόσωμα έχει μικρό fitness value καταλαμβάνει μικρότερο τμήμα και άρα είναι πιο δύσκολο να επιβιώσει.

\vspace{0.4cm}     

\(\cdot\) Η διαδικασία της διασταύρωσης υλοποιείται στη συνάρτηση \textbf{crossover.m}. Με τη διασταύρωση συνδυάζονται δύο τυχαία χρωμοσώματα, δύο γονείς και παράγονται απόγονοι (offspring). Οι κύριοι μέθοδοι διασταύρωσης είναι δύο. Στην πρώτη χρησιμοποιείται η μέση τιμή των χρωμοσωμάτων των δύο γονέων, δημιουργώντας με τον τρόπο αυτόν δύο ίδιους απογόνους. Η μέθοδος αυτή αναφέρεται ως μέθοδος της μέσης τιμής. Στη δεύτερη μέθοδο παράγονται δύο απόγονοι, καθένας από τους οποίους έχει κληρονομήσει ένα τμήμα από τον ένα γονέα και το υπόλοιπο από τον άλλο. Η διαδικασία αυτή αναφέρεται ως διακριτή διασταύρωση. Η δεύτερη μέθοδος χρησιμοποιείται συνήθως όταν η κωδικοποίηση του χρωμοσώματος έχει γίνει δυαδικά άρα, αφού έχει επιλεχθεί κωδικοποίηση με πραγματικές τιμές χρησιμοποιείται η πρώτη μέθοδος. Η συνάρτηση δηλαδή, έχει ως είσοδο δύο τυχαία χρωμοσώματα (parent1, parent2) και δημιουργεί δύο νέα ίδια, στο οποίο καθένα στοιχείο \(x_{i}, i \in [1, 17] \) είναι ίσο με \(\frac{x_{1i}+x_{2i}}{2}\). Σε κάθε γενιά, τα χρωμοσώματα χωρίζονται τυχαία σε ζεύγη και κάθε ζεύγος διασταυρώνεται με πιθανότητα \(p_{c}=0.6\). Στο τέλος, ελέγχεται αν ικανοποιούνται οι περιορισμοί του προβλήματος με τη συνάρτηση constraints και αν δεν ικανοποιούνται, η αλλαγή απορρίπτεται. Επειδή οι περιορισμοί είναι γραμμικοί, ο μέσος όρος δύο έγκυρων λύσεων είναι πάντα έγκυρη λύση.

\vspace{0.4cm}     

\(\cdot\) Η διαδικασία της μετάλλαξης υλοποιείται στη συνάρτηση \textbf{mutation.m} με τη βοήθεια των συναρτήσεων \textbf{distributeChange.m} και \textbf{adaptToChange.m}. Με τη μετάλλαξη κάποιο τυχαίο χρωμόσωμα μετασχηματίζεται ώστε να υπάρχει ποικιλία στον πληθυσμό. Συγκεκριμένα, κάθε χρωμόσωμα μεταλλάσσεται με πιθανότητα \(p_{m}=0.3\) και για τυχαία επιλεγμένο πεδίο του χρωμοσώματος
αυτού, μεταβάλλεται η ποσότητα \(x_{i}\) κατά ένα \(\Delta x_{i}\), το οποίο ακολουθεί Gaussian κατανομή \(\Delta x_{i} \sim N(0,1)\), σύμφωνα με τη σχέση \(x_{i}(k)=x_{i}(k-1)+\Delta x_{i}(k)\). Λόγω των αρχικών περιορισμών του προβλήματος, μεταλλάξεις μπορούν να γίνουν μόνο στα \(x_{1}, x_{2}, x_{3}, x_{4}\) αφού, κάθε μεταβολή που συμβαίνει σε \(x_{i}, i \in [5, 17]\) είναι αποτέλεσμα μιας μεταβολής που έγινε στους αρχικούς κλάδους \(i = 1, 2, 3, 4\). Μια μεταβολή δηλαδή π.χ. στο \(x_{4}\) προκαλεί αλλαγή στα \(x_{1}, x_{2}, x_{3}\) για να εξακολουθεί να ισχύει ο περιορισμός \(x_{1}+x_{2}+x_{3}+x_{4}=100\) και στα \(x_{9}, x_{10}\) για να εξακολουθεί να ισχύει ο περιορισμός \(x_{4}-x_{9}-x_{10}=0\). Η μεταβολή αυτή τελικά, προκαλεί αλλαγές σε όλους τους κλάδους του δικτύου. Η μετάλλαξη ενός κλάδου βέβαια, δεν έχει ως αποτέλεσμα τη μεταβολή όλων των κλάδων που αναφέρονται στον περιορισμό απαραίτητα, αρκεί ένας τουλάχιστον κλάδος να αλλάξει ώστε να εξακολουθεί να ισχύει ο περιορισμός. Για το παραπάνω παράδειγμα δηλαδή, η μεταβολή \(x_{4}+\Delta x_{4}\) μπορεί να προκαλέσει μία από τις παρακάτω αλλαγές ώστε να εξακολουθεί να ισχύει \(x_{1}+x_{2}+x_{3}+x_{4}=100\) 

1) \(x_{1}-\Delta x_{4}\) 

2) \(x_{2}-\Delta x_{4}\) 

3) \(x_{3}-\Delta x_{4}\) 

4) \(x_{1}-k \cdot \Delta x_{4}\) και \(x_{2}-(1-k) \cdot \Delta x_{4}\) μοιράζοντας έτσι τη μεταβολή σε δύο κλάδους (k τυχαίος αριθμός στο διάστημα [0, 1]),  

5) \(x_{1}-k \cdot \Delta x_{4}\) και \(x_{3}-(1-k) \cdot \Delta x_{4}\) 

6) \(x_{2}-k \cdot \Delta x_{4}\) και \(x_{3}-(1-k) \cdot \Delta x_{4}\)  

7) \(x_{1}-\frac{k_{1}}{k_{1}+k_{2}+k_{3}} \cdot \Delta x_{4}\) και \(x_{2}-\frac{k_{2}}{k_{1}+k_{2}+k_{3}} \cdot \Delta x_{4}\) και \(x_{3}-\frac{k_{3}}{k_{1}+k_{2}+k_{3}} \cdot \Delta x_{4}\) μοιράζοντας έτσι τη μεταβολή σε τρεις κλάδους (\(k_{i}\) τυχαίοι αριθμοί στο διάστημα [0, 1]). Αντίστοιχα προκύπτουν οι περιπτώσεις για τα \(x_{1}, x_{2}, x_{3}\).

Στη συνάρτηση mutation επιλέγεται τυχαία ποιο από τα \(x_{1}, x_{2}, x_{3}, x_{4}\) θα μεταλλαχθεί και με τη συνάρτηση distributeChange μοιράζεται τυχαία η μεταβολή σε τουλάχιστον ένα από τα υπόλοιπα. Όταν η μεταβολή μοιράζεται σε 2 ή παραπάνω κλάδους το ποσοστό της μεταβολής που θα πάει σε κάθε κλάδο επιλέγεται και πάλι τυχαία. Μετά την μετάλλαξη, καλείται η adaptToChange η οποία μοιράζει εκ νέου, τυχαία, τις ροές στους κλάδους \(x_{i}, i \in [5, 17]\) που επηρεάζονται από τη μεταβολή, έτσι ώστε να ικανοποιούνται όλοι οι περιορισμοί. Στο τέλος, ελέγχεται αν ικανοποιούνται οι περιορισμοί του προβλήματος με τη συνάρτηση constraints και αν δεν ικανοποιούνται, η αλλαγή απορρίπτεται.

\vspace{0.4cm}

\(\cdot\) Καθώς η επιλογή, η διασταύρωση και η μετάλλαξη είναι στοχαστικές, το καλύτερο χρωμόσωμα μιας γενιάς μπορεί να χαθεί. Γι' αυτό, χρησιμοποιείται ελιτισμός (elitism): στο τέλος κάθε γενιάς, το καλύτερο χρωμόσωμα της προηγούμενης γενιάς αντικαθιστά το χειρότερο χρωμόσωμα της νέας. Έτσι, το fitness value του καλύτερου χρωμοσώματος δεν μειώνεται ποτέ.

\vspace{0.4cm}

\(\cdot\) Η συνθήκη τερματισμού \(|\frac{f(x_{k-1})-f(x_{k})}{f(x_{k-1})}|<\epsilon\), με \(\epsilon=10^{-6}\), εφαρμόζεται στο fitness value του καλύτερου χρωμοσώματος. Επειδή ο αλγόριθμος είναι στοχαστικός και μια γενιά μπορεί να μη φέρει βελτίωση κατά τύχη, ο αλγόριθμος τερματίζει όταν η συνθήκη ισχύει για 500 διαδοχικές γενιές (ή μετά από το πολύ 20000 γενιές).

\newpage

\(\cdot\) Παρακάτω παρουσιάζεται το γράφημα του fitness value του καλύτερου χρωμοσώματος και του μέσου fitness value του πληθυσμού συναρτήσει των γενεών, καθώς και τα αποτελέσματα του αλγορίθμου, που περιέχουν τον αριθμό γενεών που χρειάστηκαν για να τερματίσει ο αλγόριθμος, το καλύτερο χρωμόσωμα, τον ελάχιστο χρόνο διάσχισης του δικτύου ανά όχημα και αν η λύση είναι έγκυρη ή όχι.

\begin{figure}[H]
    \centering
    \includegraphics[width=0.7\linewidth]{GA1_plot.jpg}
    \caption{Το fitness value συναρτήσει των γενεών για τον πρώτο αλγόριθμο}
    \label{fig:enter-label}
\end{figure}

\begin{verbatim}
The number of generations is 2004
The best chromosome is
 x1-x6:   29.94  16.72  24.03  29.31  18.56  11.38
 x7-x12:   9.86   6.87  13.41  15.91  11.20  19.86
 x13-x17: 13.24   8.34  26.14  26.90  27.10
The total travel time per vehicle is 174.8702 min
Solution is valid.
\end{verbatim}

Τα δέκα χρωμοσώματα της πρώτης γενιάς παράγονται τυχαία άρα, τα αποτελέσματα του αλγορίθμου είναι διαφορετικά σε κάθε εκτέλεση. Τα παραπάνω αποτελούν ενδεικτικό παράδειγμα. Παρατηρείται ότι με το πέρασμα των γενεών το fitness value του καλύτερου χρωμοσώματος αυξάνεται γρήγορα στις πρώτες γενιές και μετά βελτιώνεται σιγά σιγά, δηλαδή ο αλγόριθμος κατευθύνεται στη βέλτιστη λύση. Το μέσο fitness value του πληθυσμού παρουσιάζει διακυμάνσεις, αφού σε κάθε γενιά οι μεταλλάξεις δημιουργούν και χειρότερες λύσεις, οι οποίες διατηρούν την ποικιλία του πληθυσμού.

\vspace{0.4cm}

Για να ελέγξουμε την ποιότητα της λύσης, λύνουμε το ίδιο πρόβλημα και με τη συνάρτηση fmincon του Matlab. Αφού η συνάρτηση \(f(x)\) είναι κυρτή και οι περιορισμοί γραμμικοί, η fmincon βρίσκει το ολικό ελάχιστο, το οποίο είναι \(173.74\) min ανά όχημα. Σε 10 εκτελέσεις, ο γενετικός αλγόριθμος έδωσε ελάχιστο χρόνο από \(173.97\) έως \(177.41\) min ανά όχημα, δηλαδή το πολύ \(2\%\) μεγαλύτερο από το ολικό ελάχιστο.

\newpage

\section{Θέμα 3: Υλοποίηση δεύτερου γενετικού αλγορίθμου}

Στο θέμα 3, αλλάζει το V και παύει πια να θεωρείται σταθερό. Πλέον η τιμή του V που αντιπροσωπεύει τον ρυθμό εισερχόμενων οχημάτων στο δίκτυο είναι μεταβλητή και μπορεί μάλιστα να μεταβάλλεται μέχρι \(\pm 15\%\) της αρχικής του τιμής, δηλαδή \(85 \leq V \leq 115\). Άρα, το κάθε χρωμόσωμα πλέον αποτελείται από 18 στοιχεία δηλαδή  είναι ένα διάνυσμα 1x18 που περιέχει όλα τα \(x_{i}, i \in [1, 17]\) και το V ως 18η παράμετρο. Το αρχείο στο οποίο υλοποιείται ο αλγόριθμος είναι το \textbf{GeneticAlgorithm2.m}. 

\vspace{0.4cm}

Οι αλλαγές που χρειάστηκε να γίνουν είναι λίγες. Αρχικά, για την δημιουργία της
πρώτης γενιάς, γίνεται τυχαία επιλογή για το V στο διάστημα \([85, 115]\) και έπειτα γίνεται το μοίρασμα σε
όλες τις κατευθύνσεις όπως πριν (συνάρτηση \textbf{firstGeneration2.m}). Επίσης, στην συνάρτηση της μετάλλαξης (συνάρτηση \textbf{mutation2.m}) πλέον μπορεί να αλλάξει και η 18η παράμετρος και άρα να επιλεχθεί V βέλτιστο. Συγκεκριμένα, με πιθανότητα 1/2 μεταλλάσσονται οι ροές όπως στον πρώτο αλγόριθμο, για σταθερό V, ενώ αλλιώς το V μεταβάλλεται κατά \(\Delta V \sim N(0,1)\), μέσα στα όρια \([85, 115]\). Η μεταβολή \(\Delta V\) προστίθεται σε ένα τυχαίο από τα \(x_{1}, x_{2}, x_{3}, x_{4}\), ώστε να ισχύει \(x_{1}+x_{2}+x_{3}+x_{4}=V\), και η adaptToChange προσαρμόζει τα υπόλοιπα \(x_{i}\). Στη διασταύρωση δεν χρειάστηκαν ιδιαίτερες αλλαγές, μόνο αύξηση του μεγέθους του χρωμοσώματος, ενώ η συνάρτηση constraints ελέγχει τους περιορισμούς με την τιμή του V κάθε χρωμοσώματος και επιπλέον ότι \(85 \leq V \leq 115\). Τέλος, η fitness function χρησιμοποιεί το V κάθε χρωμοσώματος για τον υπολογισμό του χρόνου ανά όχημα.

\vspace{0.4cm}

\(\cdot\) Παρακάτω παρουσιάζεται το γράφημα του fitness value του καλύτερου χρωμοσώματος και του μέσου fitness value του πληθυσμού συναρτήσει των γενεών, καθώς και τα αποτελέσματα του αλγορίθμου, που περιέχουν τον αριθμό γενεών που χρειάστηκαν για να τερματίσει ο αλγόριθμος, το καλύτερο χρωμόσωμα, το βέλτιστο V (η τελευταία τιμή του χρωμοσώματος), τον ελάχιστο χρόνο διάσχισης του δικτύου ανά όχημα και αν η λύση είναι έγκυρη ή όχι.

\begin{figure}[H]
    \centering
    \includegraphics[width=0.7\linewidth]{GA2_plot.jpg}
    \caption{Το fitness value συναρτήσει των γενεών για τον δεύτερο αλγόριθμο}
    \label{fig:enter-label}
\end{figure}

\begin{verbatim}
The number of generations is 3436
The best chromosome is
 x1-x6:   27.13  13.54  19.99  24.48  16.02  11.12
 x7-x12:   8.69   4.85  11.68  12.80   9.73  18.29
 x13-x17:  8.49   6.89  21.41  22.90  22.53
V is 85.1398
The total travel time per vehicle is 123.2892 min
Solution is valid.
\end{verbatim}

Όπως στον πρώτο αλγόριθμο, τα δέκα χρωμοσώματα της πρώτης γενιάς παράγονται τυχαία άρα, τα αποτελέσματα του αλγορίθμου είναι διαφορετικά σε κάθε εκτέλεση. Τα παραπάνω αποτελούν ενδεικτικό παράδειγμα. Η τελική τιμή του V καταλήγει πάντα πολύ κοντά στο κατώτερο όριο \(V=85\) (σε 10 εκτελέσεις από \(85.00\) έως \(85.66\)), το οποίο είναι και το αναμενόμενο αποτέλεσμα. Όσο λιγότερα οχήματα εισέρχονται στο δίκτυο, τόσο λιγότερη είναι η κίνηση σε κάθε δρόμο και άρα τόσο μικρότερος ο χρόνος κίνησης \(T_{i}\), οπότε ο χρόνος διάσχισης ανά όχημα αυξάνεται με το V. Έτσι, προκύπτει μικρότερος χρόνος διάσχισης ανά όχημα σε σχέση με τον πρώτο αλγόριθμο.

\vspace{0.4cm}

Με τη συνάρτηση fmincon, για \(V=85\) το ολικό ελάχιστο είναι \(122.01\) min ανά όχημα. Σε 10 εκτελέσεις, ο δεύτερος γενετικός αλγόριθμος έδωσε ελάχιστο χρόνο από \(122.93\) έως \(131.01\) min ανά όχημα, δηλαδή το πολύ \(7\%\) μεγαλύτερο από το ολικό ελάχιστο. Η απόκλιση είναι λίγο μεγαλύτερη από τον πρώτο αλγόριθμο, αφού ο χώρος αναζήτησης είναι μεγαλύτερος λόγω της επιπλέον παραμέτρου V. Όλες οι λύσεις που προέκυψαν ικανοποιούν τους περιορισμούς του προβλήματος.
    
\end{flushleft}

\end{document}






